% Slide "How many streams the stencil needs" (new in presentation5, after the user asked how many streams the
% flat arrays of the layout slide stand for). Same strip, numbers and colours as figures/fzyx.tex: the
% 3 x 3 x 3 cube stored x fastest, then y, then z; cell 13 in the middle (blue), its 6 face neighbours (teal).
% Top: if ONE array held all 7 values, the reads at 12, 13, 14 (same row, same cache line) form one stream,
% and 10 (row y-1), 16 (row y+1), 4 (plane z-1), 22 (plane z+1) one each: 5 streams (the classic 7-point
% stencil: 5 read streams + 1 written = 6, the old bullet of the stream slide).
% Bottom: our kernel (waLBerla fzyx, pull, lbmpy order): every value has its own array and is read at one place
% only, x - c_i: f_0 at 13, f_1 (N, c = (0,1,0)) at 13 - 3 = 10, f_2 (S) at 16, f_3 (W, c = (-1,0,0)) at 14,
% f_4 (E) at 12, f_5 (T, c = (0,0,1)) at 13 - 9 = 4, f_6 (B) at 22 (checked against the generated StreamCollide,
% see figures/stencil-d3q19.tex). So one stream per array: 19 read, plus 19 written (AB pattern) and 4 outputs
% (rho, u_x, u_y, u_z) = 42 (the generated loop has 19 prefetch_read + 23 prefetch_write, one per array).
% In place (AA): reads and writes use the same 19 arrays at the same addresses: 19 + 4 = 23.
% Overlays: 1 one array, 5 streams; 2 our kernel, 7 arrays, 7 streams; 3 all 19 arrays and the total.
% presentation7-v2 (user's review: "missing the important thing: the cache line and where it falls"): redrawn
% with real distances instead of the 3 x 3 x 3 toy array. Each piece is ONE cache line of 64 B = 16 f32 values
% (cache-line.tex). Real f32 strides of our arrays (fzyx.tex): row 68 values = 272 B, plane 4488 values =
% 17 952 B (about 18 KB), array 296 208 values = 1 184 832 B (about 1.2 MB). The cell x is drawn at position 6 of
% its line, so x-1 and x+1 lie in the same line (at a line border one of them falls into the next line, which the
% loop needs next anyway).
% Overlays: 1 one array holding all 7 values: 5 streams; 2 our kernel: 7 arrays, 7 streams, + 12 more = 19.
\begin{tikzpicture}[x=1cm,y=1cm]
  \tikzset{
    vv/.style={draw=psgray!45,line width=0.2pt},
    cline/.style={draw=black!80,line width=1.0pt},
    me/.style={draw=psblue!80!black,fill=psblue!55,line width=0.4pt},
    nb/.style={draw=psteal!70!black,fill=psteal!45,line width=0.4pt},
    hd/.style={font=\scriptsize,anchor=south west,inner sep=0pt},
    tl/.style={font=\tiny,text=psgray!90!black,inner sep=1pt,align=center},
    sn/.style={circle,draw=psteal!70!black,fill=white,font=\tiny\bfseries,text=psteal!50!black,
               inner sep=0.4pt,minimum size=3.3mm},
    arrlbl/.style={font=\scriptsize,anchor=east,inner sep=1pt},
  }
  \path (0,-0.3) rectangle (8.5,4.55);   % the same size on every overlay
  \def\w{0.075}\def\h{0.3}
  %   \piece{x}{y}{cells to mark: list of pos/style}: one cache line of 16 values
  \def\piece#1#2#3{%
    \foreach \k in {0,...,15} \draw[vv] (#1+\k*\w,#2) rectangle ++(\w,\h);
    \foreach \p/\st in #3 \draw[\st] (#1+\p*\w,#2) rectangle ++(\w,\h);
    \draw[cline] (#1,#2) rectangle ++(16*\w,\h);}

  % ================= 1: one array holding all 7 values ===============================================
  \node[hd] at (0,4.32) {If \textbf{one array} held all 7 values: \textbf{5 streams}};
  \foreach \i/\px/\marks/\t/\n in {
      0/0.15/{6/nb}/{plane $z{-}1$\\18\,KB before}/4,
      1/1.9/{6/nb}/{row $y{-}1$\\272\,B before}/2,
      2/3.65/{5/nb,7/nb,6/me}/{row $y$:\\$x{-}1$, $x$, $x{+}1$}/1,
      3/5.4/{6/nb}/{row $y{+}1$\\272\,B after}/3,
      4/7.15/{6/nb}/{plane $z{+}1$\\18\,KB after}/5} {
    \piece{\px}{3.35}{\marks}
    \node[tl,anchor=north] at (\px+8*\w,3.32) {\t};
    \node[sn] at (\px+8*\w,3.92) {\n};
    \draw[psteal!70!black,line width=0.4pt] (\px+8*\w,3.68) -- (\px+8*\w,3.76);
  }
  \foreach \gx in {1.35,3.1,4.85,6.6}
    \node[font=\scriptsize,text=psgray] at (\gx+0.13,3.5) {$\cdots$};
  \node[tl,anchor=south east] at (8.5,4.3) {each box: one cache line, 16 values};

  % ================= 2: our kernel, every value in its own array =====================================
  \begin{scope}[visible on=<2->]
    \node[hd] at (0,2.42) {In our kernel, \textbf{each value has its own array}: \textbf{7 streams}};
    \foreach \a/\marks/\t/\yy in {0/{6/me}/{the cell itself}/2.0, 1/{6/nb}/{row $y{-}1$}/1.68,
        2/{6/nb}/{row $y{+}1$}/1.36, 3/{7/nb}/{$x{+}1$}/1.04, 4/{5/nb}/{$x{-}1$}/0.72,
        5/{6/nb}/{plane $z{-}1$}/0.4, 6/{6/nb}/{plane $z{+}1$}/0.08} {
      \node[arrlbl] at (1.05,\yy+0.12) {array $f_{\a}$};
      \piece{1.15}{\yy}{\marks}
      \node[tl,anchor=west] at (2.4,\yy+0.13) {\t};
      \pgfmathtruncatemacro\n{\a+1}
      \node[sn] at (3.75,\yy+0.15) {\n};
    }
    \draw[decorate,decoration={brace,amplitude=3pt},psgray!90!black] (4.05,2.3) -- (4.05,0.08)
      node[midway,right=4pt,font=\scriptsize,align=left,text width=4.1cm]
      {Each array starts 1.2\,MB after the previous one. So every read falls into a different cache line,
       and each array is a stream of its own.};
  \end{scope}
  \node[font=\scriptsize,anchor=west,inner sep=0pt,visible on=<2->] at (0,-0.2)
    {+ 12 more arrays ($f_7$ to $f_{18}$), one stream each: \textbf{19 streams read}};
\end{tikzpicture}%
